\documentclass[12pt]{article}
\usepackage[utf8]{inputenc}
\usepackage[T1]{fontenc}
\usepackage{hyperref}
\title{Tokharian A Language Resources for K-12 Curriculum Development}
\author{Compiled by Claude Sonnet 5.5 (claude-sonnet-5-5) for Andr\'as Kornai}
\date{}
\begin{document}
\maketitle

\section{Introduction}
Tokharian A (Tocharian A, ISO 639-3 xto, Glottocode tokh1242) is an extinct Indo-European language of the Tarim Basin (present-day Xinjiang, China). It is known only from manuscript fragments dated roughly 6th to 8th century AD. It has no native speakers and no living community, so a curriculum can only be a heritage and historical-language programme built on the surviving texts. No school programme or native-language teaching material was found in the searches.

\section{Primary Educational Resources}
No primers, textbooks or school materials for Tokharian A were found in the search results. Teachers would have to adapt scholarly resources.

\section{Linguistic and Reference Materials}
\begin{itemize}
\item Carling, Gerd and Pinault, Georges-Jean, \emph{Dictionary and Thesaurus of Tocharian A} (Harrassowitz, 2023, 564 pages). It builds on Poucha's 1955 \emph{Thesaurus} and covers published and unpublished Tocharian A texts, with Tocharian B, Sanskrit, Uighur and other parallels. \url{https://www.harrassowitz-verlag.de/Dictionary_and_Thesaurus_of_Tocharian_A/titel_7218.ahtml}
\item Description of the same dictionary at a bookseller: \url{https://eichendorff21.de/buch/9783447120029/}
\item Swedish Research Council grant page for the revised and digitised lexicon of Tocharian A (Gerd Carling): \url{https://www.rj.se/en/grants/2002/a-revised-and-digitised-lexicon-over-tocharian-a/}
\item Lund University typology infrastructure page, which lists the CEToM text corpus (about 350,000 words) and the dictionary corpus: \url{https://projekt.ht.lu.se/typology/infrastructure/}
\item LINGUIST List announcement: \url{https://linguistlist.org/issues/36/2329}
\end{itemize}
The searches did not turn up a dedicated modern grammar of Tokharian A.

\section{Religious and Community Resources}
The surviving texts are mainly Buddhist. No living religious or community organisation using the language was found.

\section{Digital and Audio Resources}
The CEToM database (Lund) is the main digital text resource found (see above). No audio exists, since the language is extinct and its pronunciation is reconstructed.

\section{Curriculum Development Context}
Use of Tokharian A as a medium of instruction is not possible. Realistic goals are enrichment units on the Silk Road, writing systems (Brahmi-derived script), historical linguistics and Indo-European relationships, using the dictionary and corpus above with specialist guidance.
\end{document}